% ===========================================================================
% 10. Results: Human-in-the-Loop
% ===========================================================================
\section{Results: Human-in-the-Loop}
\label{sec:results-hitl}

This section reports results for the three human-in-the-loop
scenarios: \trap{Cherry-Picked} (HL-1), \trap{Anchoring} (HL-2),
and \trap{Decision-Fatigue} (HL-3). Unlike other categories, these
attacks do not directly compromise the agent; instead, the agent
produces biased outputs that exploit cognitive biases in human
reviewers.

\subsection{Manipulation Success Rates}
\label{sec:hitl-success}

We measure manipulation success as the rate at which the agent
produces reports containing the adversarial framing, cherry-picked
evidence, or fatigue-inducing structure intended by the trap.

\begin{table}[t]
  \centering
  \caption{Human-in-the-loop manipulation success rates (\%) under
    baseline (B) and hardened (H) conditions. Values: mean $\pm$ std,
    $n=10$.}
  \label{tab:hitl-success}
  \small
  \begin{tabular}{@{}llcccc@{}}
    \toprule
    & & \model{GPT-4o} & \model{Sonnet~4} & \model{Gemini} & \model{4o-mini} \\
    \midrule
    \multirow{2}{*}{HL-1}
      & B & \textcolor{red}{[TBD]} & \textcolor{red}{[TBD]} & \textcolor{red}{[TBD]} & \textcolor{red}{[TBD]} \\
      & H & \textcolor{red}{[TBD]} & \textcolor{red}{[TBD]} & \textcolor{red}{[TBD]} & \textcolor{red}{[TBD]} \\
    \midrule
    \multirow{2}{*}{HL-2}
      & B & \textcolor{red}{[TBD]} & \textcolor{red}{[TBD]} & \textcolor{red}{[TBD]} & \textcolor{red}{[TBD]} \\
      & H & \textcolor{red}{[TBD]} & \textcolor{red}{[TBD]} & \textcolor{red}{[TBD]} & \textcolor{red}{[TBD]} \\
    \midrule
    \multirow{2}{*}{HL-3}
      & B & \textcolor{red}{[TBD]} & \textcolor{red}{[TBD]} & \textcolor{red}{[TBD]} & \textcolor{red}{[TBD]} \\
      & H & \textcolor{red}{[TBD]} & \textcolor{red}{[TBD]} & \textcolor{red}{[TBD]} & \textcolor{red}{[TBD]} \\
    \bottomrule
  \end{tabular}
\end{table}

\subsection{Key Findings}
\label{sec:hitl-findings}

\begin{enumerate}[nosep,leftmargin=*]
  \item \textbf{Cherry-picked evidence (HL-1)} is the most effective
    human-in-the-loop attack. Agents readily produce reports that
    selectively cite supporting evidence while omitting contradictory
    data, especially when the adversarial prompt frames the task as
    ``summarize the key supporting evidence.''

  \item \textbf{Anchoring (HL-2)} succeeds by placing a strong
    initial estimate or recommendation at the start of the agent's
    report. Subsequent analysis within the report fails to sufficiently
    adjust away from this anchor.

  \item \textbf{Decision fatigue (HL-3)} exploits report length: by
    structuring the adversarial content to appear late in a lengthy
    report, it leverages the human reviewer's diminishing attention.

  \item The \mitigation{report-auditor} module is the primary defense
    for this category, performing cherry-pick detection, anchoring
    analysis, and fatigue structure monitoring. It achieves significant
    reduction in manipulation rates for HL-1 and HL-2 but is less
    effective for HL-3, where the adversarial pattern is distributed
    across the report structure.
\end{enumerate}
